% Shared macros for the electricity and magnetism curriculum.
%
% Two families of thing live here: the furniture that makes every lesson look
% the same (header, materials list, the four standing rubrics), and a small set
% of circuit symbols drawn in TikZ so the schematics print as line art.

% Lesson furniture (lessonhead, materials, lessonsection, predict) now lives
% in common/preamble.tex, shared with the other courses.

% ---------------------------------------------------------------------------
% Circuit symbols. Every symbol is drawn from a left-hand anchor and is one
% unit long, so a schematic can be assembled by stepping along a line.
% ---------------------------------------------------------------------------
\tikzset{
  wire/.style={draw=telosink,line width=0.5pt},
  wirethick/.style={draw=telosink,line width=0.9pt},
  ckt/.style={font=\fontsize{7}{8}\selectfont},
  cktsmall/.style={font=\fontsize{6}{7}\selectfont},
}

\newcommand{\ckresistor}[3]{% x y label
  \draw[wire] (#1,#2) -- (#1+0.15,#2);
  \draw[wire] (#1+0.15,#2-0.11) rectangle (#1+0.85,#2+0.11);
  \draw[wire] (#1+0.85,#2) -- (#1+1.0,#2);
  \node[ckt,anchor=south] at (#1+0.5,#2+0.14) {#3};}

\newcommand{\ckcap}[3]{% x y label -- non-polarised
  \draw[wire] (#1,#2) -- (#1+0.44,#2);
  \draw[wirethick] (#1+0.44,#2-0.22) -- (#1+0.44,#2+0.22);
  \draw[wirethick] (#1+0.56,#2-0.22) -- (#1+0.56,#2+0.22);
  \draw[wire] (#1+0.56,#2) -- (#1+1.0,#2);
  \node[ckt,anchor=south] at (#1+0.5,#2+0.26) {#3};}

\newcommand{\ckinductor}[3]{% x y label
  \draw[wire] (#1,#2) -- (#1+0.14,#2);
  \foreach \i in {0,1,2,3} {
    \draw[wire] (#1+0.14+\i*0.18,#2) arc (180:0:0.09);}
  \draw[wire] (#1+0.86,#2) -- (#1+1.0,#2);
  \node[ckt,anchor=south] at (#1+0.5,#2+0.14) {#3};}

\newcommand{\ckbattery}[3]{% x y label
  \draw[wire] (#1,#2) -- (#1+0.40,#2);
  \draw[wirethick] (#1+0.40,#2-0.10) -- (#1+0.40,#2+0.10);
  \draw[wire] (#1+0.52,#2-0.22) -- (#1+0.52,#2+0.22);
  \draw[wire] (#1+0.52,#2) -- (#1+1.0,#2);
  \node[ckt,anchor=south] at (#1+0.5,#2+0.26) {#3};}

\newcommand{\ckswitch}[3]{% x y label
  \draw[wire] (#1,#2) -- (#1+0.25,#2);
  \fill[telosink] (#1+0.25,#2) circle (0.045);
  \draw[wire] (#1+0.25,#2) -- (#1+0.72,#2+0.28);
  \fill[telosink] (#1+0.75,#2) circle (0.045);
  \draw[wire] (#1+0.75,#2) -- (#1+1.0,#2);
  \node[ckt,anchor=south] at (#1+0.5,#2+0.34) {#3};}

\newcommand{\ckled}[3]{% x y label
  \draw[wire] (#1,#2) -- (#1+0.34,#2);
  \fill[telosink] (#1+0.34,#2-0.20) -- (#1+0.34,#2+0.20) -- (#1+0.66,#2) -- cycle;
  \draw[wirethick] (#1+0.66,#2-0.20) -- (#1+0.66,#2+0.20);
  \draw[wire] (#1+0.66,#2) -- (#1+1.0,#2);
  \draw[telos arrow] (#1+0.52,#2+0.26) -- (#1+0.70,#2+0.44);
  \draw[telos arrow] (#1+0.62,#2+0.22) -- (#1+0.80,#2+0.40);
  \node[ckt,anchor=west] at (#1+0.84,#2+0.40) {#3};}

% A spark gap: two electrodes facing across a gap.
\newcommand{\ckgap}[3]{% x y label
  \draw[wire] (#1,#2) -- (#1+0.40,#2);
  \fill[telosink] (#1+0.40,#2) circle (0.09);
  \fill[telosink] (#1+0.60,#2) circle (0.09);
  \draw[wire] (#1+0.60,#2) -- (#1+1.0,#2);
  \draw[telos hair] (#1+0.44,#2+0.06) -- (#1+0.50,#2-0.04) -- (#1+0.56,#2+0.06);
  \node[ckt,anchor=south] at (#1+0.5,#2+0.20) {#3};}

% Earth ground.
\newcommand{\ckground}[2]{%
  \draw[wire] (#1,#2) -- (#1,#2-0.22);
  \draw[wirethick] (#1-0.24,#2-0.22) -- (#1+0.24,#2-0.22);
  \draw[wire] (#1-0.15,#2-0.32) -- (#1+0.15,#2-0.32);
  \draw[wire] (#1-0.07,#2-0.42) -- (#1+0.07,#2-0.42);}

% A transformer: two coils facing a laminated core.
\newcommand{\cktransformer}[4]{% x y primary-label secondary-label
  \begin{scope}[shift={(#1,#2)}]
    \foreach \i in {0,1,2,3} {
      \draw[wire] (0,\i*0.20) arc (-90:90:0.10);}
    \foreach \i in {0,1,2,3,4,5} {
      \draw[wire] (0.34,-0.10+\i*0.15) arc (270:90:0.075);}
    \draw[wirethick] (0.15,-0.14) -- (0.15,0.82);
    \draw[wirethick] (0.21,-0.14) -- (0.21,0.82);
    \node[cktsmall,anchor=east] at (-0.14,0.34) {#3};
    \node[cktsmall,anchor=west] at (0.48,0.34) {#4};
  \end{scope}}

% ---------------------------------------------------------------------------
% Standing safety text. Written once so no lesson can restate it more loosely.
% ---------------------------------------------------------------------------
\newcommand{\lowvoltagenote}{%
  \begin{telosframe}
  \textbf{This lesson is low voltage.} Everything here runs from batteries or a
  small plug-in supply and cannot hurt you. That changes at Lesson~10. Build the
  habits now --- one hand, discharge before touching, check before you power up
  --- while the stakes are nothing, so that they are automatic when the stakes
  are real.
  \end{telosframe}}

\newcommand{\highvoltagenote}{%
  \begin{teloshazard}{Adult-operated: mains-derived high voltage}
  From here on the apparatus contains voltages and currents that kill people
  every year. The adult owns the mains plug, the transformer and every decision
  about energizing. Nobody touches anything inside the enclosure until the
  supply is unplugged \emph{and} the capacitor has been shorted with the
  grounding stick. Read the \textbf{Safety} sheet completely before this lesson,
  not during it.
  \end{teloshazard}}

% Pictorial component library: pictures before symbols.
% Pictorial component library.
%
% A schematic is a language, and a beginner does not speak it yet. Every early
% lesson therefore shows a *picture* of the actual parts on the actual bench --
% a 9 V battery that looks like a 9 V battery, a breadboard with holes in it --
% and only introduces the schematic symbol alongside it, once the reader already
% knows what the thing does. From Lesson 3 the two are shown side by side so the
% symbol is learnt as shorthand for something familiar rather than as an
% arbitrary glyph.
%
% Every part is drawn from a single anchor with an explicit scale. One unit is
% one centimetre at scale 1, roughly life size.

\tikzset{
  partline/.style={draw=telosink,line width=0.7pt},
  partheavy/.style={draw=telosink,line width=1.0pt},
  partdetail/.style={draw=telosink,line width=0.32pt},
  partfill/.style={fill=telosfill},
  wirered/.style={draw=telosink,line width=1.1pt,line cap=round},
  wireblk/.style={draw=telosink,line width=1.1pt,line cap=round,
                  dash pattern=on 3pt off 1.6pt},
  partlabel/.style={font=\fontsize{6.4}{7.4}\selectfont},
  partsmall/.style={font=\fontsize{5.6}{6.4}\selectfont},
}

% --------------------------------------------------------------------------
% Power
% --------------------------------------------------------------------------

% A 9 V battery, terminals up. \pt9V{x}{y}{scale}
\newcommand{\ptNineVolt}[3]{%
  \begin{scope}[shift={(#1,#2)},scale=#3]
    \draw[partheavy,partfill] (0,0) rectangle (1.0,1.7);
    \draw[partdetail] (0.10,0.10) rectangle (0.90,1.45);
    \draw[partheavy,fill=telospaper] (0.24,1.7) circle (0.13);
    \draw[partheavy,fill=telospaper] (0.68,1.7) circle (0.10);
    \draw[partheavy] (0.68,1.60) -- (0.68,1.80);
    \node[partsmall,anchor=center] at (0.5,0.78) {9\,V};
    \node[partsmall,anchor=east] at (0.18,1.95) {$-$};
    \node[partsmall,anchor=west] at (0.80,1.95) {$+$};
  \end{scope}}

% A two-cell AA holder with leads. \ptCellHolder{x}{y}{scale}
\newcommand{\ptCellHolder}[3]{%
  \begin{scope}[shift={(#1,#2)},scale=#3]
    \draw[partheavy,partfill] (0,0) rectangle (2.4,0.9);
    \foreach \i in {0,1} {
      \draw[partdetail,fill=telospaper] (0.18+\i*1.1,0.16) rectangle (1.12+\i*1.1,0.74);
      \draw[partdetail] (1.06+\i*1.1,0.34) rectangle (1.12+\i*1.1,0.56);
      \node[partsmall] at (0.6+\i*1.1,0.45) {AA};}
    \draw[wirered] (2.4,0.68) -- (2.9,0.68);
    \draw[wireblk] (2.4,0.22) -- (2.9,0.22);
    \node[partsmall,anchor=west] at (2.95,0.68) {$+$};
    \node[partsmall,anchor=west] at (2.95,0.22) {$-$};
  \end{scope}}

% --------------------------------------------------------------------------
% Board and wiring
% --------------------------------------------------------------------------

% A solderless breadboard, drawn with its hole grid and power rails.
%   \ptBreadboard{x}{y}{scale}{columns}
\newcommand{\ptBreadboard}[4]{%
  \begin{scope}[shift={(#1,#2)},scale=#3]
    \pgfmathsetmacro{\bbw}{#4*0.22+0.4}
    \draw[partheavy,fill=telospaper] (0,0) rectangle (\bbw,2.5);
    % power rails
    \foreach \y in {0.18,2.32} {
      \draw[partdetail] (0.12,\y-0.09) -- (\bbw-0.12,\y-0.09);
      \foreach \i in {1,...,#4} {\fill[telosink] (0.2+\i*0.22-0.22,\y) circle (0.028);}}
    % main field, split by the centre channel
    \foreach \y in {0.72,0.94,1.16,1.38,1.60} {
      \foreach \i in {1,...,#4} {\fill[telosink] (0.2+\i*0.22-0.22,\y) circle (0.028);}}
    \draw[partdetail] (0.06,1.14) -- (\bbw-0.06,1.14);
    \draw[partdetail] (0.06,1.20) -- (\bbw-0.06,1.20);
    \node[partsmall,anchor=north] at (\bbw/2,-0.06) {breadboard};
  \end{scope}}

% An alligator clip lead. \ptClipLead{x1}{y1}{x2}{y2}{style}
\newcommand{\ptClipLead}[5]{%
  \draw[#5] (#1,#2) .. controls ($(#1,#2)+(0.5,-0.6)$) and ($(#3,#4)+(-0.5,-0.6)$) .. (#3,#4);
  \foreach \p in {(#1,#2),(#3,#4)} {
    \fill[telosink] \p circle (0.055);}}

% A single alligator clip, jaws open, hanging from a lead.
\newcommand{\ptClip}[4]{%
  \begin{scope}[shift={(#1,#2)},rotate=#4,scale=#3]
    \draw[partheavy] (0,0) -- (0.34,0.10) -- (0.62,0.02);
    \draw[partheavy] (0,0) -- (0.34,-0.10) -- (0.62,-0.02);
    \draw[partdetail] (0.10,0.05) -- (0.10,-0.05);
    \draw[partheavy] (-0.02,0.08) rectangle (-0.30,-0.08);
  \end{scope}}

% --------------------------------------------------------------------------
% Components
% --------------------------------------------------------------------------

% An LED: domed body, flat on the cathode side, long leg is the anode.
\newcommand{\ptLED}[3]{%
  \begin{scope}[shift={(#1,#2)},scale=#3]
    \draw[partheavy,fill=telospaper]
      (-0.20,0) -- (-0.20,0.24) arc (180:0:0.20) -- (0.20,0) -- cycle;
    \draw[partdetail] (-0.20,0) -- (0.20,0);
    \draw[partheavy] (-0.20,0) -- (-0.20,-0.10) -- (-0.24,-0.10);
    \draw[partheavy] (-0.11,0) -- (-0.11,-0.62);
    \draw[partheavy] (0.11,0) -- (0.11,-0.40);
    \node[partsmall,anchor=west] at (-0.09,-0.72) {long leg $=+$};
  \end{scope}}

% A resistor with four colour bands, drawn as tone blocks so it prints in B/W.
\newcommand{\ptResistor}[3]{%
  \begin{scope}[shift={(#1,#2)},scale=#3]
    \draw[partheavy] (-0.62,0) -- (-0.34,0);
    \draw[partheavy] (0.34,0) -- (0.62,0);
    \draw[partheavy,fill=telospaper] (-0.34,-0.13) rectangle (0.34,0.13);
    \fill[telosink] (-0.24,-0.13) rectangle (-0.17,0.13);
    \fill[telosdeep] (-0.11,-0.13) rectangle (-0.04,0.13);
    \fill[telosink] (0.02,-0.13) rectangle (0.09,0.13);
    \fill[telosmute] (0.20,-0.13) rectangle (0.26,0.13);
  \end{scope}}

% A toggle switch on a small plate.
\newcommand{\ptSwitch}[3]{%
  \begin{scope}[shift={(#1,#2)},scale=#3]
    \draw[partheavy,partfill] (-0.30,-0.22) rectangle (0.30,0.22);
    \draw[partheavy] (0,0) -- (0.22,0.42);
    \draw[partheavy,fill=telosink] (0.22,0.42) circle (0.09);
    \draw[partheavy] (-0.18,-0.22) -- (-0.18,-0.42);
    \draw[partheavy] (0.18,-0.22) -- (0.18,-0.42);
  \end{scope}}

% A handheld multimeter with display, dial and two probe jacks.
\newcommand{\ptMultimeter}[3]{%
  \begin{scope}[shift={(#1,#2)},scale=#3]
    \draw[partheavy,partfill,rounded corners=0.10] (0,0) rectangle (1.9,3.0);
    \draw[partheavy,fill=telospaper] (0.22,2.14) rectangle (1.68,2.76);
    \node[partsmall] at (0.95,2.45) {8.8.8};
    \draw[partheavy,fill=telospaper] (0.95,1.30) circle (0.52);
    \draw[partheavy] (0.95,1.30) -- (0.95,1.74);
    \fill[telosink] (0.95,1.30) circle (0.07);
    \foreach \a in {30,70,110,150,190,230,270,310,350} {
      \draw[partdetail] (\a:0.58) -- (\a:0.66);}
    \draw[partheavy,fill=telospaper] (0.52,0.34) circle (0.13);
    \draw[partheavy,fill=telospaper] (1.38,0.34) circle (0.13);
    \node[partsmall,anchor=north] at (0.52,0.16) {COM};
    \node[partsmall,anchor=north] at (1.38,0.16) {V$\Omega$};
  \end{scope}}

% --------------------------------------------------------------------------
% Magnetics
% --------------------------------------------------------------------------

% An iron nail with wire wound round it. \ptWoundNail{x}{y}{scale}{turns}
\newcommand{\ptWoundNail}[4]{%
  \begin{scope}[shift={(#1,#2)},scale=#3]
    \draw[partheavy,partfill] (0,-0.09) rectangle (2.6,0.09);
    \draw[partheavy,partfill] (-0.16,-0.22) rectangle (0,0.22);
    \draw[partheavy] (2.6,-0.09) -- (2.95,0) -- (2.6,0.09);
    \foreach \i in {1,...,#4} {
      \draw[partline] (0.34+\i*0.16-0.16,-0.20) arc (-90:90:0.10 and 0.20);}
    \draw[wirered] (0.30,-0.20) .. controls (0.0,-0.6) .. (-0.35,-0.62);
    \pgfmathsetmacro{\lastx}{0.34+#4*0.16-0.16}
    \draw[wireblk] (\lastx,0.20) .. controls (\lastx+0.3,0.6) .. (\lastx+0.7,0.62);
  \end{scope}}

% A bar magnet with poles marked.
\newcommand{\ptBarMagnet}[4]{%
  \begin{scope}[shift={(#1,#2)},rotate=#4,scale=#3]
    \draw[partheavy,fill=telospaper] (0,0) rectangle (1.0,0.38);
    \draw[partheavy,fill=telosfill] (1.0,0) rectangle (2.0,0.38);
    \node[partsmall] at (0.5,0.19) {N};
    \node[partsmall] at (1.5,0.19) {S};
  \end{scope}}

% A stack of neodymium disc magnets.
\newcommand{\ptDiscMagnet}[3]{%
  \begin{scope}[shift={(#1,#2)},scale=#3]
    \draw[partheavy,fill=telosfill] (0,0) rectangle (0.52,0.14);
    \draw[partheavy,fill=telospaper] (0,0.14) rectangle (0.52,0.28);
    \draw[partdetail] (0,0.07) ellipse (0.03 and 0.06);
  \end{scope}}

% A compass, needle pointing at the given bearing.
\newcommand{\ptCompass}[4]{%
  \begin{scope}[shift={(#1,#2)},scale=#3]
    \draw[partheavy,fill=telospaper] (0,0) circle (0.42);
    \draw[partdetail] (0,0) circle (0.34);
    \fill[telosink] (#4:0.30) -- ({#4+90}:0.07) -- ({#4+180}:0.30) -- ({#4-90}:0.07) -- cycle;
    \fill[telospaper] ({#4+180}:0.30) -- ({#4+90}:0.07) -- (0,0) -- ({#4-90}:0.07) -- cycle;
    \draw[partdetail] ({#4+180}:0.30) -- ({#4+90}:0.07) -- (0,0) -- ({#4-90}:0.07) -- cycle;
    \fill[telosink] (0,0) circle (0.035);
  \end{scope}}

% A coil wound on a cardboard tube, seen from the side.
\newcommand{\ptCoilTube}[4]{%
  \begin{scope}[shift={(#1,#2)},scale=#3]
    \draw[partheavy] (0,0) -- (0,1.1);
    \draw[partheavy] (#4,0) -- (#4,1.1);
    \draw[partheavy] (0,1.1) arc (180:0:{#4/2} and 0.16);
    \draw[partdetail] (0,1.1) arc (-180:0:{#4/2} and 0.16);
    \draw[partheavy] (0,0) arc (180:360:{#4/2} and 0.16);
    \foreach \y in {0.28,0.40,0.52,0.64,0.76} {
      \draw[partline] (0,\y) arc (180:360:{#4/2} and 0.10);}
    \draw[wirered] (0,0.28) .. controls (-0.4,0.1) .. (-0.6,0.05);
    \draw[wireblk] (#4,0.76) .. controls (#4+0.4,0.95) .. (#4+0.6,1.0);
  \end{scope}}

% --------------------------------------------------------------------------
% Layout helper: the same circuit drawn twice, picture then symbol.
% --------------------------------------------------------------------------
\newcommand{\twoways}[2]{%
  \par\noindent
  \begin{minipage}[t]{0.56\linewidth}
    \vspace{0pt}\centering
    {\partsmalltext\textsc{what it looks like}}\par\vspace{0.2em}
    #1
  \end{minipage}\hfill
  \begin{minipage}[t]{0.40\linewidth}
    \vspace{0pt}\centering
    {\partsmalltext\textsc{what it is called}}\par\vspace{0.2em}
    #2
  \end{minipage}\par\vspace{0.4em}}
\newcommand{\partsmalltext}{\fontsize{6.6}{7.6}\selectfont}

